\documentclass{situationist-preprint}

\title{Sophie's Choice Reconsidered}
\msid{ETH-15-b}
\msfamily{ETH-15}
\msdiscipline{Normative and Applied Ethics}
\mstype{Companion manuscript}
\mskeywords{moral dilemmas; Sophie's choice; tie-breaking; picking; remorse; infringement}

\begin{document}
\maketitle

\begin{abstract}
Sophie's choice is the case most often offered as a genuine moral dilemma: a mother forced to choose which of her two children will be killed, on pain of both being killed. ETH-15-a argues that there are no genuine dilemmas and that every apparent dilemma has exactly one permissible option. We work through the case. If the children's claims are exactly balanced, the uniquely permissible option is to select without preference; if one child has the better prospect of survival, it is to choose that child, though the mother may be able to identify it only in retrospect. In neither version does she act wrongly. Her remorse is nonetheless fitting, because remorse responds to the infringement of a claim, not to wrongdoing.
\end{abstract}

\section{The case}

In Styron's novel, a mother arriving at Auschwitz with her two children is told by an officer that she may keep one of them; the other will be sent to its death, and if she does not choose, both will be \citep{Styron1979}. Three options are available: to choose the first child, to choose the second, and to refuse. Each infringes a requirement in force, the requirement of a parent to protect each child, so the case is an apparent dilemma in the sense of Definition~2.3 of ETH-15-a. By Theorem~3.4 of ETH-15-a, it has exactly one permissible option. We ask which.

\paragraph{Refusal.} Refusing leads to the death of both children. It infringes both requirements and satisfies neither, and it is worse on the balance of reasons than either choice. By Assumption~2.5 of ETH-15-a, it is not permissible. The thought that refusal keeps the mother's hands clean mistakes the source of the requirement: her duty is to protect her children, not to avoid involvement in the officer's crime.

\paragraph{Exact balance.} Suppose the children's claims are exactly balanced: each has the same prospect of surviving the camp, and nothing else distinguishes them in a way that bears on permissibility. Then, by Definition~2.1 of ETH-15-a, the options include not only taking one child or the other directly but also selecting between them without preference, by picking in the sense of \citet{UllmannMargalitMorgenbesser1977}. By Assumption~2.7 of ETH-15-a, an impartial selection is better than either partial choice, since it gives each child's claim equal standing. It is therefore the uniquely permissible option.

\paragraph{Unequal prospects.} Suppose instead that one child has the better prospect of surviving the camp. Then the requirements are not balanced: choosing that child gives the greater chance that one child survives, and it is better on the balance of reasons. By Assumption~2.5 of ETH-15-a, it is uniquely permissible. Which child has the better prospect is a fact about what the camp will bring, and the mother cannot know it when she chooses. The uniquely permissible option exists at the time of choice, but she may be able to identify it only in retrospect, as Section~4 of ETH-15-a explains.

\section{The mother does not act wrongly}

\begin{proposition}\label{p:notwrong}
In neither version does the mother act wrongly in choosing one child.
\end{proposition}

\begin{proof}
In the version with exact balance, a selection made under terror and in seconds, without a ranking of the children, is a picking: it is not the expression of a preference for one child but the only way left of keeping one alive. It is therefore the impartial option, which is permissible. In the version with unequal prospects, either she chooses the child with the better prospect, which is permissible, or she chooses the other on her best evidence, in which case she takes the impermissible option without fault. By Principle~P1 of ETH-13-a, her culpability depends only on her will and evidence, and so she is not to blame; and what she does is no failure of regard for either child, but a failure to know a fact that no one could know. The judgement that she acted wrongly, which presents her act as a failure to do what she ought on her evidence and in her power, is mistaken in both versions.
\end{proof}

In the second version, the proposition is carefully stated. If the mother chooses, on her evidence, the child whose prospect turns out to be worse, her option was not the permissible one, and this is a fact about what she did. But it is not a wrong for which she is answerable as a wrongdoer. ETH-13 calls this the scope of responsibility as distinct from culpability: the outcome falls within what she did, and she is not more culpable for it.

\section{Remorse}

The mother's remorse is the datum that defenders of dilemmas take to show that she did wrong whatever she chose \citep{Williams1965}. ETH-15-a explains the datum without a dilemma: the permissible option still infringes a requirement in force (Corollary~3.5 of ETH-15-a). ETH-27-a supplies the account of the emotion. Remorse presents its object as an act that infringed, through one's own agency, a serious claim of another (Definition~2.3 of ETH-27-a), and an emotion is fitting if and only if things are as it presents them (Principle~P1 of ETH-27-a). The child the mother did not choose had a claim on her protection, and her choice infringed it. Her remorse is therefore fitting. Guilt, understood as presenting her act as wrong, is not fitting, by the reasoning of Corollary~3.3 of ETH-27-a; nor is the blame of others.

\begin{objection}
To call the mother's remorse fitting while denying that she did wrong is to tell her that she has something to be remorseful for and nothing to answer for. The combination is unintelligible.
\end{objection}

\begin{reply}
She has something to answer for: she infringed her child's claim, and the duties of acknowledgment that ETH-27-a derives (Proposition~3.4 of ETH-27-a) apply to her as to any agent who infringes a serious claim in acting rightly. What she does not have is a violation, an infringement that was impermissible. The wrong in the case was done, and the blame belongs, to the officer who made her choice an instrument of his cruelty.
\end{reply}

\section{Conclusion}

Sophie's choice has a uniquely permissible option in each of its versions: the unpreferring choice under exact balance, and the choice of the child more likely to survive under unequal prospects, identifiable perhaps only in retrospect. The mother does not act wrongly in choosing, and her remorse is fitting because she infringed a claim, not because she violated one.

\begin{thebibliography}{9}\small
\bibitem[Styron(1979)]{Styron1979} Styron, W. (1979). \emph{Sophie's Choice}. New York: Random House.
\bibitem[Ullmann-Margalit and Morgenbesser(1977)]{UllmannMargalitMorgenbesser1977} Ullmann-Margalit, E. and Morgenbesser, S. (1977). Picking and choosing. \emph{Social Research}, 44(4), 757--785.
\bibitem[Williams(1965)]{Williams1965} Williams, B. (1965). Ethical consistency. \emph{Proceedings of the Aristotelian Society, Supplementary Volumes}, 39, 103--124.
\end{thebibliography}

\end{document}
